\documentclass[../main.tex]{subfiles}\begin{document}
\section{Frontend}
\begin{figure}[H]
    \centering
    \includegraphics[width=1\linewidth]{IMG/BILAG/TEKNISKANALYSE/frontend-mock.png}
    \caption{Mock frontend}
    \label{fig:mock-frontend}
\end{figure}

Følgende beskriver bla bla bla.\\
Det er internt i gruppen blevet valgt og udvikle en Web app, men at den skal kunne understøtte tilgang igennem en mobil. Dette valg er taget, da gruppen har generel mere erfaring med udvikling af Web apps fremfor Mobile apps
Der er lavet en analyse på baggrund af følgende valg: Angular, React og Vue.

\subsection{Krav}
Hvad har \textbf{jeres produkt} brug for?\\
Eksemple: Backend-teknologien skal understøtte REST API'er, authentication, relationel databaseadgang og asynkron behandling. Da projektets udviklingsperiode er begrænset, prioriteres development speed og eksisterende kompetencer højt.

\subsection*{Kandidater}
Eksemple: ASP.NET Core, Go og Node.js.

\subsection*{Kriterier}
Definer kriterier for teknologien

\subsection*{Analyse}
Diskuter teknologiernes faktiske styrker og svagheder.

\subsection{Scoring}
\begin{tabular}{|l|l|l|l|l|}
    \hline
    \rowcolor{gray!30}
    \textbf{Kriterie} & \textbf{Vægt} & \textbf{Tek 1} & \textbf{Tek 2} & \textbf{Tek 3}\\
    \hline
    Projekt fit & 25\% & - & - & - \\
    \hline
    Udviklingshastighed & 20\% & - & - & - \\
    \hline
    Maintainability & 15\% & - & - & - \\
    \hline
    Performance & 10\% & - & - & - \\
    \hline
    Ecosystem & 10\% & - & - & - \\
    \hline
    Team kompetence & 20\% & - & - & - \\
    \hline
    \rowcolor{gray!30}
    Vægtet scorer & 100\% & - & - & - \\
    \hline
\end{tabular}

\subsection*{Beslutning}
\end{document}